% ------------------------------------------------------------------ cover
\thispagestyle{blank}
\begin{tikzpicture}[remember picture,overlay]
  \fill[disc] (current page.south west) rectangle (current page.north east);
  % the 8088's twenty address lines, A19 to A0
  \addressbandlabels{(current page.north west)+(0,{-3.05in-\bleed})}{\paperwidth}{0.42in}{white!70!disc}
  \node[anchor=north,text=white!70!disc,inner sep=0pt,font=\sffamily\small] at ($(current page.north)+(0,{-3.47in-\bleed})$) {the 8088's twenty address lines: one megabyte};
  % The Lyon Mountain Company's lion, alone, small, upper right
  \node[anchor=north east,inner sep=0pt] at ($(current page.north east)+({-0.8in-\bleed},{-0.55in-\bleed})$)
    {\includegraphics[height=1.0in]{figures/lmc-lion.png}};
  \node[anchor=north west,text=white,inner sep=0pt,align=left] at ($(current page.north west)+({0.85in+\bleed},{-3.95in-\bleed})$)
    {\cond\bfseries\fontsize{68}{68}\selectfont RetrOS/LT};
  \node[anchor=north west,text=white,inner sep=0pt,align=left,font=\cond\fontsize{18}{22}\selectfont] at ($(current page.north west)+({0.9in+\bleed},{-5.0in-\bleed})$)
    {\subtitlelines};
  \node[anchor=north west,text=white,inner sep=0pt,align=left] at ($(current page.north west)+({0.9in+\bleed},{-5.95in-\bleed})$)
    {\sffamily\bfseries\fontsize{24}{28}\selectfont \booktitle};
  \node[anchor=north west,text=white!80!disc,inner sep=0pt,align=left,font=\sffamily\fontsize{11.5}{15}\selectfont] at ($(current page.north west)+({0.9in+\bleed},{-6.5in-\bleed})$)
    {RetrOS/LT Version 1.0, DOS 2.10\\for the IBM Personal Computer, the XT, and compatibles};
  % the processor
  \node[anchor=south,inner sep=0pt] at ($(current page.south)+(0,{1.25in+\bleed})$)
    {\dipchip{3.1}};
\end{tikzpicture}
\newpage
% Two-sided: the inside of the front cover is left empty.
\ifmwtwoside\thispagestyle{blank}\null\newpage\fi

% ------------------------------------------------------------------ title page
\frontpages
\setcounter{page}{1}
\pagestyle{blank}
\vspace*{0.4in}
\hspace*{-\dimexpr\marginparwidth+\marginparsep\relax}%
\begin{minipage}[t]{\dimexpr\textwidth+\marginparwidth+\marginparsep\relax}
\raggedright
{\cond\bfseries\fontsize{40}{44}\selectfont\color{disc} RetrOS/LT\par}
\vspace{6pt}
{\cond\fontsize{17}{21}\selectfont\color{disc} \subtitleline\par}
\vspace{14pt}
\begin{tikzpicture}\addressband{(0,0)}{\linewidth}{4pt}\end{tikzpicture}\par
\vspace{14pt}
{\sffamily\bfseries\fontsize{22}{26}\selectfont \booktitle\par}
\vspace{8pt}
{\large RetrOS/LT Version 1.0, for the IBM Personal Computer,\\
the Personal Computer XT, and compatible computers\par}
\vspace{0.5in}
\hspace*{0.3in}\dipchip{2.0}\par
\vspace{24pt}
{\small\textbf{The Lyon Mountain Company}\\
Republic Lyon Mountain Corporation\\
Lyon Mountain, New York 12952 USA\par}
\end{minipage}
\newpage

% ------------------------------------------------------------------ edition notice
\pagestyle{blank}
\vspace*{\fill}
\begin{wide}
\small
\textbf{\edition}

This edition applies to Version 1.0 of RetrOS/LT, build \release, and to
all later releases until otherwise indicated in new editions. Changes are
made periodically to the information herein. \editionchanges

A form for readers' comments is at the back of this book. If the form has
been removed, comments may be addressed to P.O. Box 290, Lyon Mountain,
New York 12952 USA.

\textbf{Where RetrOS/LT comes from.} RetrOS/LT is derived from the MS-DOS 2
source code that Microsoft Corporation released under the MIT License.
That code is Copyright \copyright\ Microsoft Corporation; the notice that
must accompany it is printed in Appendix~D and is on the RetrOS/LT
diskettes. RetrOS/LT is not MS-DOS and is not PC DOS. It is an independent
product of The Lyon Mountain Company, and is not affiliated with, endorsed
by, or supported by Microsoft Corporation or International Business
Machines Corporation.

\textbf{RetrOS/LT is licensed, not sold}, under the RetrOS/LT License
printed in Appendix~D: a Single-Machine License, for one computer at a
time, or a Blanket License, for all the computers one owner owns or
leases.

\textbf{Trademarks.} IBM, Personal Computer XT and PC DOS are trademarks
of International Business Machines Corporation. Microsoft and MS-DOS are
trademarks of Microsoft Corporation. Intel is a trademark of Intel
Corporation. Other names used in this book may be trademarks of their
owners.

\copyright\ Copyright The Lyon Mountain Company, Republic Lyon Mountain Corporation 2026
\end{wide}
\newpage

% ------------------------------------------------------------------ preface
\pagestyle{plain}
\renewcommand\runningtitle{About This Book}
\chapter*{About This Book}
\addcontentsline{toc}{chapter}{About This Book}

This book explains how RetrOS/LT is made: what it takes from Microsoft's
MS-DOS 2 source and what it adds, how the computer starts, how the
commands IBM wrote for PC DOS 2.10 were written again, how Install and the
diskettes work, and how all of it is built and tested. It is the companion
to the \textit{RetrOS/LT Guide to Operations}, and is written for two
readers at once.

\begin{plainwords}
If you do not write programs, read these panels. Each one tells, without
any code, what the part it opens does, and why it matters to the person
using the computer.
\end{plainwords}

If you do write programs, the rest is for you: the source, a listing at a
time, with the reasons behind it. Most listings are in the assembly
language of the Intel 8088, as DOS itself is.

\section{How This Book Is Organized}

\begin{tabularx}{\textwidth}{@{}l>{\raggedright\arraybackslash}X@{}}
\toprule
\textbf{Section} & \textbf{About}\\
\midrule
1 Overview & what RetrOS/LT is made of, and where each part comes from\\
2 Building & the source as published, the 1988 tools in an emulator, the patches, NASM, and what comes out\\
3 Starting the Computer & from the master boot record to the prompt: the startup record and \file{RETBIO.COM}\\
4 The Kernel & Microsoft's code as RetrOS/LT builds it, and the console functions the published source lacks\\
5 The IBM Commands & the commands and driver written again to behave as IBM's\\
6 Install & \file{INSTALL.COM}, and the two diskettes\\
7 Testing & QEMU, MAME's IBM PC 5150, and the review that followed\\
Appendixes & the source files, the patches, and credits\\
\bottomrule
\end{tabularx}

\section{Listings}

RetrOS/LT keeps its changes to Microsoft's files as a series of patches,
not as edited files, so a file such as \file{SYSINIT.ASM} exists in
RetrOS/LT's form only while it is being built. This book's listings are
taken the same way: when the book is made, \file{scripts/prep.py} prepares
the source as the build does, and the lines are taken from the result. A
listing of Microsoft's code therefore shows the file as RetrOS/LT builds
it, numbered as the file then is. Listings of RetrOS/LT's own files are
taken from the repository. If the code a listing points at changes, the
book will not build until the listing, and the text about it, are
checked.

\section{Conventions}

Numbers ending in \code{h} are hexadecimal: \code{0FF0h} is 4,080, and
\code{0070:0000} is a segment and offset. \code{AX}, \code{CS} and the
like are the processor's registers. Names of routines in the source are
printed \routine{LIKETHIS}, and the patches by their numbers,
\patch{0002}. Notes for programmers are headed \textbf{Engineering note}:

\engnote{An engineering note explains a choice that is not obvious from
the code, often one made after an earlier attempt was wrong.}

\mwnewpart
\renewcommand\runningtitle{Contents}
\renewcommand\contentsname{Contents}
\tableofcontents
